\input{configTD}

% ── Poly à trous : commandes spécifiques ──────────────────────────

\newcommand{\atrou}[2][3cm]{%
  \ifthenelse{\boolean{showSolutions}}{%
    \par\vspace{0.5em}%
    \begin{mdframed}[backgroundcolor=ThemeLight, linewidth=0pt,
      skipabove=0pt, skipbelow=0pt,
      innertopmargin=8pt, innerbottommargin=8pt]%
    #2%
    \end{mdframed}%
    \par\vspace{0.5em}%
  }{%
    \par\vspace{0.5em}%
    \noindent\fbox{\begin{minipage}[c][#1]%
      {\dimexpr\textwidth-2\fboxsep-2\fboxrule\relax}%
      \color{white}\rule{0pt}{0pt}%
    \end{minipage}}%
    \par\vspace{0.5em}%
  }%
}

\newcommand{\val}[1]{%
  \ifthenelse{\boolean{showSolutions}}{#1}{\phantom{#1}}%
}

\newcommand{\esquisse}[1][5.5cm]{%
  \begin{center}%
  \ifthenelse{\boolean{showSolutions}}{%
    \fbox{\begin{minipage}[c][#1]{0.85\textwidth}%
      \centering\vfill
      \textit{\small Croquis réalisé en cours.}%
      \vfill
    \end{minipage}}%
  }{%
    \fbox{\begin{minipage}[c][#1]{0.85\textwidth}%
      \hfill\vfill\hfill
    \end{minipage}}%
  }%
  \end{center}%
}

\newcommand{\degC}{\text{\textdegree C}}

\title{\sffamily\bfseries Analyse Numérique --- CM3\\[0.3em]
  \Large Problème inverse, méthode de Newton et ajustement de paramètres}
\author{}
\date{23 avril 2026}

\begin{document}
\sffamily
\maketitle
\thispagestyle{empty}

% ====================================================================
\section*{1. Le problème --- une dalle de béton sur un chantier}
% ====================================================================

\begin{center}
\begin{mdframed}[backgroundcolor=ThemeLight, linecolor=Theme,
  linewidth=1.5pt, innertopmargin=10pt, innerbottommargin=10pt]
\centering\large\itshape
Sur un chantier, une dalle de béton a été coulée à $45\,\degC$.\\
Un capteur mesure sa température à 60 minutes : $24{,}03\,\degC$.\\
On peut marcher sur la dalle si sa température est inférieure à $18\,\degC$ ?\\
\textbf{Dans combien de temps pourra-t-on y circuler\,?}
\end{mdframed}
\end{center}

\bigskip

\textbf{Données.}
\atrou[2cm]{%
\begin{itemize}[leftmargin=*]
  \item Température initiale : $T_0 = 45\,\degC$.
  \item Température ambiante : $T_{\mathrm{amb}} = 15\,\degC$.
  \item Mesure du capteur à $t = 60$ min : $T_{\mathrm{mes}} = 24{,}03\,\degC$.
  \item Critère de sécurité : circulation possible dès que $T < 18\,\degC$.
\end{itemize}
}
\medskip

\textbf{Modèle.}
\atrou[2cm]{%
Comme au CM1 pour le café :
\[
  \frac{dT}{dt} = -k\,(T - T_{\mathrm{amb}}),
  \qquad T(0) = T_0.
\]}

\medskip

\textbf{Question.}
Pour prédire quand $T(t) = 18$, de quoi a-t-on besoin\,?

\atrou[2cm]{%
Du coefficient d'échange thermique~$k$.
Il dépend du chantier (vent, humidité, bâche...).
On ne peut pas le mesurer directement.
}

\medskip

\textbf{Ce qu'on sait faire (CM1, TD4).}
Si $k$ est connu, on sait calculer $T(t)$ :
soit par la formule exacte (si on peut résoudre),
soit par un schéma numérique
(Euler, Crank-Nicolson, RK4).

\medskip

\textbf{On veut faire l'inverse :}
partir d'une mesure
et retrouver le paramètre $k$.

% ====================================================================
\section*{2. Reformulation --- annuler le résidu}
% ====================================================================

Pour un $k$ donné, on peut calculer la température prédite
à $t = 60$ min :
\atrou[2cm]{%
\[
  T_{\mathrm{pred}}(k) = 15 + 30\,e^{-60\,k}.
\]}

Le bon $k$ est celui qui fait coïncider prédiction et mesure :
\[
  T_{\mathrm{pred}}(k) = T_{\mathrm{mes}}
\]

\medskip

\textbf{La fonction résidu:}
\[
  g(k) = T_{\mathrm{pred}}(k) - T_{\mathrm{mes}} = 15 + 30\,e^{-60\,k} - 24{,}03.
\]

On cherche $k$ tel que : 
\atrou[1cm]{%
\[
  g(k) = 0.
\]
}

\medskip

\textbf{Question.}
Peut-on résoudre cette équation en $k$\,?

\atrou[10cm]{%
elle est exponentielle en $k$.

On pourrait l'inverser à la main dans ce cas précis
grâce au logarithme, mais en pratique $T_{\mathrm{pred}}(k)$ est souvent calculée
par un schéma numérique et aucune formule explicite n'est disponible.

\smallskip

Il nous faut donc une méthode générale
pour résoudre $g(k) = 0$.

\[
k = 
\]
}

\bigskip

\begin{center}
\begin{mdframed}[backgroundcolor=ThemeLight, linecolor=Theme,
  linewidth=1pt, innertopmargin=8pt, innerbottommargin=8pt]
\centering\itshape
Comment résoudre $g(x) = 0$
quand $g$ est non linéaire
et qu'on ne sait pas inverser sa formule\,?
\end{mdframed}
\end{center}
  
\newpage
% ====================================================================
\section*{3. Cas bâtiment (problème) --- local technique}
% ====================================================================

\begin{center}
\begin{mdframed}[backgroundcolor=ThemeLight, linecolor=Theme,
  linewidth=1.5pt, innertopmargin=10pt, innerbottommargin=10pt]
\centering\large\itshape
Local technique d'un bâtiment tertiaire.\\
Le chauffage à 22 degrés est coupé à 20h.\\
À 22h, le capteur intérieur mesure $19\,\degC$.\\
\textbf{Peut-on identifier le coefficient global de pertes thermiques $U$ ?}
\end{mdframed}
\end{center}

\medskip

\textbf{Modèle simplifié.}
\[
  C\,\frac{dT_{\mathrm{int}}}{dt}
  = -U\,\bigl(T_{\mathrm{int}} - T_{\mathrm{ext}}(t)\bigr),
  \qquad T_{\mathrm{int}}(20h)=22\,\degC.
\]
Le paramètre inconnu est $U$ (en W/K).

\medskip

\textbf{On fixe la capacité thermique globale $C=4{,}0\times10^5\;\mathrm{J/K}$ ?}

elle mesure l'énergie nécessaire pour faire varier
la température intérieure de $1$ K.



\medskip

\textbf{Données extérieures (mesures météo).}
\[
  T_{\mathrm{ext}}(20h)=12,\;
  T_{\mathrm{ext}}(21h)=10,\;
  T_{\mathrm{ext}}(22h)=8,\;
\]

\textbf{Question.}
Comment définir $T_{\mathrm{ext}}(t)$ entre ces heures ?

\atrou[5cm]{%
On interpole les mesures (interpolation linéaire, TD1)
pour obtenir une fonction continue
$T_{\mathrm{ext}}(t)$ sur $[20h,22h]$
\[
  T_{\mathrm{ext}}(t) = 12 - 2\,(t - 20).
\]
}

On peut donc résoudre l'équation : 
\atrou[5cm]{%
\[
  T(t) = K e^{(-U/C)(t - 20)} -2t + 52+2C/U
\]
On détermine ensuite $K$ en utilisant la condition initiale :
\[
  T(20h) = 22.
\]
\[
  22 = K e^{(-U/C)(20 - 20)} -2*20 + 52+2C/U
\]
\[
  22 = K -40 + 52+2C/U
\]
\[
  K = 22 - 12 - 2C/U
  \]
}
\medskip

\textbf{Problème inverse.}
On définit le résidu :
\[
  g(U) = T_{\mathrm{int,pred}}(23h;U)-19.
  g(U) = K(U) e^{(-U/C)(23 - 20)} -2*23 + 52+2C/U - 19
\]
On cherche $g(U)=0$.

\medskip

\begin{center}
\begin{mdframed}[backgroundcolor=ThemeLight, linecolor=Theme,
  linewidth=1pt, innertopmargin=8pt, innerbottommargin=8pt]
\centering\itshape
Ce cas n'est pas résoluble à la main :\\
Comment faire  ?
\end{mdframed}
\end{center}
% ====================================================================
\section*{4. Méthode de Newton}
% ====================================================================

\textbf{Idée.}
On ne sait pas résoudre $g(x) = 0$ en général,
mais on sait résoudre un problème \emph{linéaire}. \newline
À chaque étape, on \textbf{remplace $g$ par sa tangente}
et on cherche le zéro de cette tangente.

\medskip

\esquisse[5.5cm]

\textbf{Construction.}
Partons d'un $x_n$ proche de la racine cherchée.
La tangente à $g$ en $x_n$ a pour équation :
\[
  y = g(x_n) + g'(x_n)\,(x - x_n).
\]

\medskip

\textbf{Question.}
Pour quelle valeur de $x$ cette tangente
coupe-t-elle l'axe des abscisses\,?

\atrou[2.5cm]{%
On résout $g(x_n) + g'(x_n)\,(x - x_n) = 0$ :
\[
  x = x_n - \frac{g(x_n)}{g'(x_n)}.
\]
On prend ce nouveau $x$ comme point
de départ de l'itération suivante.
}

\medskip

\textbf{La méthode de Newton.}
À partir d'un $x_0$ donné, on itère :
\[
  \boxed{\;x_{n+1} = x_n - \frac{g(x_n)}{g'(x_n)}.\;}
\]

\medskip

\textbf{Vocabulaire.}
\begin{itemize}[leftmargin=*]
  \item La courbe a été \textbf{linéarisée}
        (on remplace $g$ par sa tangente).
  \item À chaque itération, on résout
        le \textbf{problème linéaire} et pas le problème exact.
  \item On recommence à partir de la nouvelle valeur.
\end{itemize}

% ── Exemple scalaire d'appropriation ──────────────────────────────

\medskip

\textbf{Exemple pour s'approprier la méthode.}
Calculons $\sqrt{2}$ à la main.
On cherche $x > 0$ tel que $x^2 = 2$, soit
\[
  f(x) = x^2 - 2 = 0.
\]

\textbf{Question.}
Écrire la formule de Newton pour ce $f$.

\atrou[2cm]{%
$f'(x) = 2x$, donc :
\[
  x_{n+1} = x_n - \frac{x_n^2 - 2}{2\,x_n}
  = \frac{1}{2}\left(x_n + \frac{2}{x_n}\right).
\]
\textit{(On retrouve la célèbre \og méthode de Héron\fg{}
pour calculer une racine carrée.)}
}

\medskip

\textbf{Itérations à partir de $x_0 = 2$} :

\atrou[5.5cm]{%
\begin{center}
\renewcommand{\arraystretch}{1.3}
\begin{tabular}{c|c|c|c}
  $n$ & $x_n$ & $f(x_n) = x_n^2 - 2$ & Chiffres corrects \\
  \hline
  $0$ & $2$                   & $2$                    & $0$ \\
  $1$ & $\tfrac12(2 + 1) = 1{,}5$        & $0{,}25$              & $1$ \\
  $2$ & $\tfrac12(1{,}5 + 4/3) \approx 1{,}41\overline{6}$ & $0{,}006\ldots$       & $3$ \\
  $3$ & $\approx 1{,}414215\ldots$       & $2 \times 10^{-6}$    & $6$ \\
  $4$ & $\approx 1{,}41421356\ldots$     & $\approx 10^{-12}$    & $12$ \\
\end{tabular}
\end{center}

\smallskip

Valeur exacte : $\sqrt{2} \approx 1{,}41421356$.
}

\medskip

\textbf{Observation.}
Le nombre de chiffres corrects \textbf{double}
à chaque itération.

\medskip

\begin{center}
\begin{mdframed}[backgroundcolor=ThemeLight, linecolor=Theme,
  linewidth=1pt, innertopmargin=8pt, innerbottommargin=8pt]
\centering\itshape
\textbf{Convergence quadratique.}\\
Si $x_0$ est assez proche de la racine et si $g'$ ne s'annule pas,
l'erreur se comporte comme $\varepsilon_{n+1} \sim C\,\varepsilon_n^2$ :\\
le nombre de chiffres corrects double à chaque itération.
\end{mdframed}
\end{center}

\medskip

\textbf{Conditions d'application --- à retenir.}

\atrou[3.5cm]{%
\begin{itemize}[leftmargin=*]
  \item Il faut pouvoir évaluer $g(x)$ \emph{et} $g'(x)$.
  \item Le point de départ $x_0$ doit être
        \textbf{pas trop loin} de la racine.
  \item $g'(x)$ ne doit pas s'annuler
        près de la racine (sinon division par zéro,
        ou instabilité numérique).
  \item En échange de ces contraintes :
        convergence très rapide (quadratique).
\end{itemize}
}

% ====================================================================
\section*{5. Application de Newton --- local technique}
% ====================================================================

On applique maintenant Newton au \textbf{deuxième problème}
(local technique), posé en section~3.

\medskip

\textbf{Rappel du résidu.}
\[
  g(U)=T_{\mathrm{int,pred}}(23h;U)-19.
\]

\textbf{Mise à jour de Newton.}
\[
  U_{n+1}=U_n-\frac{g(U_n)}{g'(U_n)},
  \qquad
  g'(U_n)\approx\frac{g(U_n+\delta U)-g(U_n)}{\delta U},
  \quad \delta U=5.
\]

\medskip

\textbf{Point de départ.}
A priori, on ne connaît pas $U$.
On prend une estimation raisonnable : $U_0=90\;\mathrm{W/K}$.

\medskip

\textbf{Pas 1 --- fait au tableau (sorties RK4 données).}

\atrou[5cm]{%
\begin{itemize}[leftmargin=*]
  \item RK4 donne $T_{\mathrm{int,pred}}(23h;90)=20{,}4$,
        donc $g(U_0)=20{,}4-19=\mathbf{+1{,}4}$.
  \item RK4 donne $T_{\mathrm{int,pred}}(23h;95)=20{,}1$,
        donc $g(95)=\mathbf{+1{,}1}$.
  \item \[
          g'(U_0)\approx\frac{1{,}1-1{,}4}{5}=-0{,}06.
        \]
  \item \[
          U_1=90-\frac{1{,}4}{-0{,}06}\approx\mathbf{113{,}3}.
        \]
\end{itemize}
}

\medskip

\textbf{Pas 2 --- à vous.}
On vous donne :
\[
  T_{\mathrm{int,pred}}(23h;113{,}3)=19{,}2,\qquad
  T_{\mathrm{int,pred}}(23h;118{,}3)=18{,}9.
\]
Calculer $g(U_1)$, $g'(U_1)$, puis $U_2$.

\atrou[5.5cm]{%
\begin{itemize}[leftmargin=*]
  \item $g(U_1)=19{,}2-19=\mathbf{+0{,}2}$.
  \item $g(118{,}3)=18{,}9-19=\mathbf{-0{,}1}$.
  \item \[
          g'(U_1)\approx\frac{-0{,}1-0{,}2}{5}=-0{,}06.
        \]
  \item \[
          U_2=113{,}3-\frac{0{,}2}{-0{,}06}\approx\mathbf{116{,}6}.
        \]
\end{itemize}
}

\medskip

\textbf{Pas 3 (donné).}

\smallskip

Une itération de plus donne
$U_3 \approx 117\;\mathrm{W/K}$.

\medskip

\textbf{Bilan des itérations.}

\begin{center}
\renewcommand{\arraystretch}{1.3}
\begin{tabular}{c|c|c|c}
  Itération & $U_n$             & $g(U_n)$       & Chiffres corrects \\
  \hline
  $0$ & $90{,}0$             & $+1{,}4$       & $0$ \\
  $1$ & $113{,}3$            & $+0{,}2$       & $1$ \\
  $2$ & $116{,}6$            & proche de $0$  & $2$--$3$ \\
  $3$ & $117{,}0$            & $\approx 0$    & $\geq 4$ \\
\end{tabular}
\end{center}

\medskip

\textit{Le nombre de chiffres corrects double
à chaque itération : c'est la signature
de la convergence quadratique.}

% ── Prédiction ────────────────────────────────────────────────────

\medskip

\textbf{Interprétation physique.}
Le coefficient identifié est :
\[
  \boxed{U \approx 117\;\mathrm{W/K}.}
\]

\atrou[2.5cm]{%
On a calibré les pertes thermiques globales
du local sur la plage 20h--23h.
Cette valeur de $U$ permet ensuite de faire
des prévisions (durée de maintien en température,
puissance de chauffage à prévoir, etc.).
}

\medskip

\begin{center}
\begin{mdframed}[backgroundcolor=ThemeLight, linecolor=Theme,
  linewidth=1pt, innertopmargin=8pt, innerbottommargin=8pt]
\centering\itshape
On a enchaîné : interpolation des données météo
$\to$ EDO (résolue numériquement) $\to$ problème inverse
$\to$ Newton.\\
\textbf{Trois itérations} suffisent pour identifier $U$
avec une bonne précision.
\end{mdframed}
\end{center}

% ── Discussion ────────────────────────────────────────────────────

\medskip

\textbf{Deux points de recul.}

\begin{enumerate}[leftmargin=*]
  \item \textbf{Problème direct vs inverse.}
        Calculer $T$ sachant $U$ = \emph{une} résolution d'EDO.
        Retrouver $U$ sachant $T$ = \emph{plusieurs} résolutions
        d'EDO (une par itération de Newton).
        L'identification est structurellement plus coûteuse.
  \item \textbf{Et si on ne connaît pas la formule de $g'$\,?}
        On l'approche par différence finie (CM2) :
        \[
          g'(U) \approx \frac{g(U + \delta U) - g(U)}{\delta U}.
        \]
        C'est ce que fait Python en arrière-plan
        quand la formule explicite n'est pas disponible.
\end{enumerate}

% ====================================================================
\section*{6. Ouverture --- plusieurs mesures et moindres carrés}
% ====================================================================

\textbf{Jeu complet de mesures.}

\begin{center}
\renewcommand{\arraystretch}{1.3}
\begin{tabular}{|c|ccccc|}
  \hline
  $t$ (min) & $0$ & $30$ & $60$ & $90$ & $120$ \\
  \hline
  $T$ mesuré (\degC) & $45$ & $31{,}5$ & $24{,}0$ & $20{,}0$ & $17{,}7$ \\
  \hline
\end{tabular}
\end{center}

\medskip

\textbf{Question.}
Avec une seule mesure, on avait une équation $g(k) = 0$.
Avec cinq mesures, on a cinq équations en un seul inconnu.
Que faire\,?

\atrou[3cm]{%
En général, ces cinq équations
\textbf{n'ont pas de solution commune}
(bruit de mesure, modèle imparfait).

\smallskip

On ne cherche plus à les annuler toutes,
on minimise la \textbf{somme des carrés des résidus} :
\[
  \boxed{\;J(k) = \sum_{i=1}^N
    \bigl(T_{\mathrm{pred}}(t_i; k)
    - T_{\mathrm{mes},\,i}\bigr)^2.\;}
\]

On est passé d'une \emph{équation} à une \emph{optimisation}.
C'est la méthode des \textbf{moindres carrés}.
}

\medskip

\textbf{Pont avec Newton.}

\atrou[3cm]{%
Minimiser $J(k)$, c'est annuler sa dérivée :
\[
  J'(k) = 0.
\]
On est ramené à une équation non linéaire,
que l'on sait traiter par Newton :
\[
  k_{n+1} = k_n - \frac{J'(k_n)}{J''(k_n)}.
\]

\smallskip

\begin{center}
\begin{tabular}{ll}
  \hline
  \textbf{Une mesure} & Newton sur $g(k) = 0$ \\
  \textbf{Plusieurs mesures} & Newton sur $J'(k) = 0$ \\
  \hline
\end{tabular}
\end{center}
}

\medskip

\textbf{Détail du calcul.}

Notons $r_i(k) = T_{\mathrm{pred}}(t_i; k) - T_{\mathrm{mes},i}$ le résidu de la $i$-ème mesure.
On a :
\[
  J(k) = \sum_{i=1}^N r_i(k)^2.
\]

\atrou[4cm]{%
En dérivant :
\[
  J'(k) = 2\sum_{i=1}^N r_i(k)\,r_i'(k),
  \qquad
  J''(k) = 2\sum_{i=1}^N \Bigl( r_i'(k)^2 + r_i(k)\,r_i''(k) \Bigr).
\]
La condition $J'(k) = 0$ s'écrit :
\[
  \sum_{i=1}^N r_i(k)\,r_i'(k) = 0.
\]
C'est une équation non linéaire en $k$ qu'on résout par Newton.
}

\medskip

\textbf{Simplification de Gauss-Newton.}

\atrou[4cm]{%
Près de la solution, les résidus $r_i(k)$ sont petits (le modèle colle aux données).
On néglige le terme $r_i(k)\,r_i''(k)$ dans $J''$ :
\[
  J''(k) \approx 2\sum_{i=1}^N r_i'(k)^2.
\]
C'est l'approximation de \textbf{Gauss-Newton} :
elle évite de calculer les dérivées secondes $r_i''$.

\smallskip

La mise à jour devient :
\[
  k_{n+1} = k_n - \frac{\displaystyle\sum_{i=1}^N r_i(k_n)\,r_i'(k_n)}
                       {\displaystyle\sum_{i=1}^N r_i'(k_n)^2}.
\]
}

\medskip

\textbf{Illustration.}
Allure de $J(k)$ et de sa dérivée $J'(k)$ :

\esquisse[4.5cm]

\textbf{Vocabulaire culturel (hors programme).}

\begin{itemize}[leftmargin=*]
  \item \textbf{Quasi-Newton (BFGS)} : on approche $J''$
        par différences finies successives.
        Utile quand la dérivée seconde est coûteuse
        ou indisponible.
  \item \textbf{Levenberg--Marquardt} : variante spécifique
        aux moindres carrés non linéaires.
        Standard en ingénierie pour l'ajustement de modèles.
  \item En Python : \verb|scipy.optimize.curve_fit|,
        \verb|scipy.optimize.least_squares|.
\end{itemize}

\end{document}
